\documentclass{article}

% --- Language and Font Setup ---
\usepackage[english]{babel}
\usepackage{fontspec}

% Standard font for Overleaf/XeLaTeX
\setmainfont{FreeSerif} 

% --- Math and Logic ---
\usepackage{amsmath, amssymb, amsthm}
\usepackage{listings}
\usepackage{xcolor}

% --- LOOP Language Definition ---
\lstdefinelanguage{LOOP}{
    keywords={LOOP, DO, END, IF, THEN, ELSE, :=},
    sensitive=true,
    morecomment=[l]{//},
    morestring=[b]"
}

\lstset{
    language=LOOP,
    basicstyle=\ttfamily\small,
    keywordstyle=\color{blue}\bfseries,
    commentstyle=\color{gray}\itshape,
    frame=single,
    tabsize=2,
    breaklines=true,
    numbers=left,
    numberstyle=\tiny\color{gray}
}

% --- Document Info ---
\title{Computational Models: Encoding Exercises}
\author{Vasilis Ioannou}
\date{\today}

\begin{document}

\maketitle

\section*{Primitive Recursion}

\subsection*{Definition}

The class $\mathcal{PR}$ of \textbf{primitive recursive functions} is defined as the smallest class of functions $f: \mathbb{N}^k \to \mathbb{N}$ that satisfies the following conditions:

\subsection*{1. Initial (Basic) Functions}
\begin{itemize}
    \item \textbf{The Zero Function:} $Z(x) = 0$.
    \item \textbf{The Successor Function:} $S(x) = x + 1$.
    \item \textbf{The Predecessor Function:} $P(x + 1) = x$ with $P(0)=0$.
    \item \textbf{The Projection Functions:} $U_i^n(x_1, \dots, x_n) = x_i$.
\end{itemize}

\subsection*{2. Closure under Operations}
\subsubsection*{A. Composition}
$h(\mathbf{x}) = f(g_1(\mathbf{x}), \dots, g_m(\mathbf{x}))$.
\subsubsection*{B. Primitive Recursion}
$f(\mathbf{x}, 0) = g(\mathbf{x})$ and $f(\mathbf{x}, y+1) = h(\mathbf{x}, y, f(\mathbf{x}, y))$.

\pagebreak
\section*{Exercise 1}

\subsection*{Proposition: The function $\text{add}(x,y)$ is primitive recursive.} 

\paragraph{1. The Base Case ($y=0$):}
$g(x) = U_1^1(x) = x$. Thus, $\text{add}(x,0) = x$.

\paragraph{2. The Inductive Step ($y \to y+1$):}
$h(x, y, z) = S(U_3^3(x, y, z)) = z + 1$.
Thus, $\text{add}(x, y+1) = \text{add}(x, y) + 1$.

\paragraph{3. Conclusion:}
Since $g$ and $h$ are built from basic functions, $\text{add}(x,y)$ is primitive recursive.

\subsection*{Proposition: The function $\text{mult}(x,y)$ is primitive recursive.} 

\paragraph{1. The Base Case ($y=0$):}
$g(x) = Z(x) = 0$. Thus, $\text{mult}(x,0) = 0$.

\paragraph{2. The Inductive Step ($y\to y+1$):}
We use the previously defined addition function:
$h(x, y, z) = \text{add}(U_3^3(x, y, z), U_1^3(x, y, z)) = z + x$.
Thus, $\text{mult}(x, y+1) = \text{mult}(x, y) + x$.

\paragraph{3. Conclusion:}
As a composition of primitive recursive functions ($\text{add}$ and projections), $\text{mult}(x,y)$ is primitive recursive.

\subsection*{Proposition: The function $\text{div}(x,y)$ is primitive recursive.}

\paragraph{1. Helper Functions:}
We utilize the following functions:
\begin{itemize}
    \item \textbf{Proper Subtraction:$x\dot{-}y$.}
    \item \textbf{Cosignum:} $\overline{\text{sg}}(x)$, which is $1$ if $x=0$ and $0$ otherwise.
\end{itemize}

\paragraph{2. Construction via Bounded Sum:}
The value of $\text{div}(x, y)$ counts how many $i \in \{1, \dots, x\}$ satisfy $i \cdot y \leq x$:
\[
\text{div}(x, y) = \sum_{i=1}^{x} \overline{\text{sg}}((i \cdot y) \dot{-} x)
\]

\paragraph{3. Formal Verification:}
The product, proper subtraction, and cosignum are all primitive recursive. Since the class $\mathcal{PR}$ is closed under \textbf{Bounded Summation}, $\text{div}(x, y)$ is primitive recursive.

\section*{Excersise 3}

\subsection{Problem Definition}
Let $\mathcal{P}$ be the set of primitive recursive functions. Show that if $g(x, y) \in \mathcal{P}$, then the function $f(x, z)$ defined by the bounded sum:
\[
f(x, z) = \sum_{y=0}^{z} g(x, y)
\]
is also in $\mathcal{P}$.

\section*{Proof}

To prove that $f \in \mathcal{P}$, we must show that it can be defined using the basic primitive recursive functions and the operations of composition and primitive recursion.

\subsection*{1. Defining the Recursion}
We define $f(x, z)$ by induction on $z$:

\begin{itemize}
    \item \textbf{Base Case ($z=0$):}
    \[
    f(x, 0) = \sum_{y=0}^{0} g(x, y) = g(x, 0)
    \]
    
    \item \textbf{Recursive Step ($z+1$):}
    \begin{align*}
    f(x, z+1) &= \sum_{y=0}^{z+1} g(x, y) \\
    &= \left( \sum_{y=0}^{z} g(x, y) \right) + g(x, z+1) \\
    &= f(x, z) + g(x, z+1)
    \end{align*}
\end{itemize}

\subsection*{2. Mapping to the Primitive Recursion Schema}
The standard schema for primitive recursion for a function of $n+1$ variables is:
\begin{enumerate}
    \item $f(\vec{x}, 0) = h_1(\vec{x})$
    \item $f(\vec{x}, z+1) = h_2(\vec{x}, z, f(\vec{x}, z))$
\end{enumerate}

In our specific case, the parameters are:
\begin{itemize}
    \item \textbf{$h_1(x) = g(x, 0)$}: Since $g$ is primitive recursive and $0$ is a constant function, $h_1$ is PR by composition.
    \item \textbf{$h_2(x, z, w) = w + g(x, z+1)$}:
    \begin{itemize}
        \item The addition function $+$ is known to be PR.
        \item $w$ is the projection $P_3^3(x, z, w)$, which is a basic PR function.
        \item $g(x, z+1)$ is the composition of $g$ with the successor function $S(z)$.
    \end{itemize}
\end{itemize}

\subsection*{Conclusion}
Since $h_1$ and $h_2$ are both primitive recursive, and $f$ is defined via the primitive recursion schema from $h_1$ and $h_2$, it follows that $f \in \mathcal{P}$. Thus, the set of primitive recursive functions is closed under bounded summation.

\qed



\end{document}
